\section{Introduction and results at a glance}\label{sec:intro}
For a cutoff \(X\), let \(\PP(X)\) index primary degree-one Eisenstein
primes and consider the complex matrix with entries
\((\rho/\sigma)_6\), extended by zero along the diagonal.
Multiplying this matrix by a Gauss-sum vector twisted by
\(\kappa\) produces an arithmetic analogue of a structured spectral
probe. We study the quotient between the output norm and the
random-phase scale. The central observation is a \emph{measured},
not proved, crossover as the dimensionless parameter
\(t=\Nm_6(\kappa)/X\) varies.

\fig{visual_abstract}{.96}{Schematic overview of the arithmetic
object and statistic. This panel is a diagram, not experimental
evidence; formulas use the conventions of CAND-002 in
\texttt{research/CONJECTURE\_REGISTER.md} at source commit
\texttt{30965467}.}{visual}
\fig{crossover}{.98}{The primary empirical profile. Exponents
\(1\)--\(4\) and the frozen RS-S016 discovery-bin means are displayed
against \(\log_{10}t\); transparent points are finite computed classes
from RS-S015, RS-S016 and RS-S021. Sources:
\texttt{results/rs\_s015\_class\_tests\_2026-10-26.json},
\texttt{rs\_s016\_x\_dependence\_2026-10-26.json} and
\texttt{rs\_s021\_large\_x\_2026-10-26.json}, at
\texttt{30965467}. No extrapolation beyond the measured points is made.}{collapse}

The candidate comprises three distinct \emph{conjectures}:
a compact-\(t\) crossover, a large-\(t\) fifth-power deficit,
and a separate fixed-class asymptotic. The third is \textbf{not}
a consequence of the first two. Small positive \(t\) and the boundary
\(t\to0\) are different regimes of quantification. The figure is not
a proof of any limiting statement.

The finite mathematical core consists of exact symbol and Gauss-sum
relations formalized in Lean, and paper-level identities reducing the
strength statistic to a signed two-point kernel. A proof of the observed
shape would require estimates for moving-shift prime-modulus first
moments and their \emph{signed} correlations, not a fit to a figure.

The work relates to residue-matrix characterizations
\cite{ddk2016}, metaplectic Gauss-sum theory
\cite{patterson1987,kazhdan1984,broeker2013,chinta2003},
the cubic and quartic first-moment problems
\cite{dunn2024,ddhl2023}, and recent results on large-sieve
nonorthogonality and sextic norms
\cite{dfdh2026,xusang2026}. These citations provide context, not
a claim of priority or an assertion that the present asymptotics
are established.
